% ---- pg-274 （书页 262） ----

联合置信区间, 并用以考察主效应的本质.

\noindent 6.15\quad 参看例 6.13.

\noindent (a) 取 $\alpha=0.05$, 对因子 1 和因子 2 的效应作近似的 $\chi^2$(似然比)检验, 并与例 6.13 中精确 $F$ 检验的结论作比较. 解释两者的差别.

\noindent (b) 对三个响应中的每一对响应, 试用\hyperref[eq:6-70]{式(6-70)}构造因子 1 的诸效应参数之差的 95\% 联合置信区间. 对因子 2 的效应参数也作同样的计算.

\noindent\mvbold{下面的练习需要使用计算机.}

\noindent 6.16\quad 对 30 块木板中的每一块的硬度指标的四次测量结果列于\hyperref[tab:4.3]{表 4.3} (参见例 4.14). 对一块给定的木板, 测量可重复进行. 假设硬度的测量值由 4 个处理产生, 用重复测量设计中的方法检验处理的相等性(取 $\alpha=0.05$), 并对均值水平的对照构造 95\% 联合置信区间, 此处均值水平代表动态测量值与静态测量值的比较.

\noindent 6.17\quad \hyperref[tab:6.8]{表 6.8} 中的数据是从检验两个数字认知的心理学模型中收集到的. 数字的认知过程是否依赖于数字的给出方式(单词, 阿拉伯数字)?

\begin{mvtable}
\phantomsection\label{tab:6.8}\mvtabcap{表 6.8\quad 数字奇偶性数据(中值反应时间单位为毫秒)}
\begin{tabular}{rrrr}
\toprule
WordDiff & WordSame & ArabicDiff & ArabicSame\\
$(x_1)$ & $(x_2)$ & $(x_3)$ & $(x_4)$\\\midrule
869.0 & 860.5 & 691.0 & 601.0\\
995.0 & 875.0 & 678.0 & 659.0\\
1 056.0 & 930.5 & 833.0 & 826.0\\
1 126.0 & 954.0 & 888.0 & 728.0\\
1 044.0 & 909.0 & 865.0 & 839.0\\
925.0 & 856.5 & 1 059.5 & 797.0\\
1 172.5 & 896.5 & 926.0 & 766.0\\
1 408.5 & 1 311.0 & 854.0 & 986.0\\
1 028.0 & 887.0 & 915.0 & 735.0\\
1 011.0 & 863.0 & 761.0 & 657.0\\
726.0 & 674.0 & 663.0 & 583.0\\
982.0 & 894.0 & 831.0 & 640.0\\
1 225.0 & 1 179.0 & 1 037.0 & 905.5\\
731.0 & 662.0 & 662.5 & 624.0\\
975.5 & 872.5 & 814.0 & 735.0\\
1 130.5 & 811.0 & 843.0 & 657.0\\
945.0 & 909.0 & 867.5 & 754.0\\
747.0 & 752.5 & 777.0 & 687.5\\
656.5 & 659.5 & 572.0 & 539.0\\
919.0 & 833.0 & 752.0 & 611.0\\
751.0 & 744.0 & 683.0 & 553.0\\
774.0 & 735.0 & 671.0 & 612.0\\
941.0 & 931.0 & 901.5 & 700.0\\
751.0 & 785.0 & 789.0 & 735.0\\
767.0 & 737.5 & 724.0 & 639.0\\
\bottomrule
\end{tabular}
\end{mvtable}
